\documentclass{notebook}
\title{時間bin平均を考慮したモデル光度曲線の計算および観測データを用いた計算パラメータのカイ2乗フィッティング}
\author{KUMADA Ryota}
\date{\today}

\begin{document}
\maketitle

\section{動機と方針}
観測データのfluxに関する$\chi^2$フィッティングモデルは
有意水準$\alpha = 0.05$の$\chi^2$検定で大部分が棄却された。

観測データの時間方向の不確かさをちゃんと考慮することが
必要であると考えた。観測データの時間方向は\qty{1}{\min}の
binningで切られており、これはfluxのそれとは不確かさの性質が異なる。

以前までの$\chi^2$フィッティングでは、
観測データの時刻に対応するモデルfluxの値をそのまま採用していた。
これを変更する。すなわち、モデル計算で得られた光度曲線を
観測データと同じ\qty{1}{\min}のbinに切り、各binで
fluxを時間平均する。このfluxに基づいた$\chi^2$フィッティングを行う。


\section{手順}
計算で得られた光度曲線の\lstinline|DataFrame| について、\qty{60}{\s}の
商でグルーピングする。
そのためには、計算で得た光度曲線はいったん保存しておいたほうがよい？

\begin{equation}
  \ab<F_{\nu}> = \frac{1}{t_1 - t_0}\int_{t_0}^{t_1} F_{\nu}(t)\odif{t}
\end{equation}
\end{document}
